\abstract{Raman and X-ray diffraction (XRD) spectroscopies provide complementary vibrational and crystallographic fingerprints critical for inorganic material characterization. Despite the success of machine learning in materials informatics, the field lacks comprehensive, open benchmark datasets mapping raw spectral files to chemical families and material properties. We introduce Spec2Prop-InorgBench, a curated open-data Raman--XRD benchmark comprising 4,118 usable Raman spectra, 4,027 clean inorganic samples, and 1,844 property-matched samples, including 1,462 Raman--XRD linked pairs. To bridge the gap between raw instrumental output and machine learning, we develop a runtime preprocessing pipeline that transforms raw spectra into a standardized 2048-point vector, followed by chemistry-aware feature extraction yielding 32 global-shape and 126 expert spectral descriptors. We propose an AI-assisted screening workflow, Spec2Prop-Edge, which employs a LightGBM-DART classifier to predict nine major inorganic chemical families (e.g., silicates, oxides, carbonates) with an accuracy of 67.22\% and a macro-F1 score of 56.57\%. We extend this approach to spectra-to-property prediction, achieving macro-F1 scores of 61.09\% for band-gap class, 67.56\% for metal/non-metal classification, and 78.91\% for formation-energy class. Spec2Prop-Edge is designed as a confidence-aware screening assistant for candidate prioritization and preliminary material-family interpretation, rather than replacing final experimental confirmation.}

\keywords{Inorganic materials, Raman spectroscopy, X-ray diffraction, Materials informatics, Spectral preprocessing, LightGBM, Spectra-to-property prediction}
